% !TeX root = ../papers/nc-p3.tex
\section{十五次元閉包における整域窓の一様な全次数判定}
\label{sec:finite-domain-specialization-parameters}
\begin{theorem}\label{thm:finite-domain-specialization-parameters}
十五次元成分の閉包で有限四辺の標準有限Frobenius候補を取る。次は同値である。
\begin{enumerate}
\item 窓は標準Koszul左右AS正則で、一つの頂点の隅が整域である。
\item 経路長48の四つの中心元が定理\ref{thm:successively-central-cutoff}の有限検査を満たす。
すなわち中心剰余を189次、元の経路次元を564次、標準複体を192次まで調べればよい。
\end{enumerate}
このとき実際に全ての合成可能な非零経路元の積は非零であり、全中心の各次元は
十五次元成分の一般中心のものに一致する。中心は有限生成で、窓はその上有限である。
四次式が四つのパラメータに取れることは仮定しない。有限検査を外れる点に
整域でない正則窓が存在するかは、この同値だけからは決まらない。
\end{theorem}
\begin{proof}[一隅の被約traceの持ち上げ]
正則な点へ一般窓から離散付値環で特殊化する。全経路次数で一般点と特殊点の次元が一致し、
各次数の経路格子は自由である。$A=e_0Be_0$ を整域と仮定した隅とし、次数を4で割る。
定理\ref{thm:quantum-component-intersection}の証明により、一般の隅の可除代数には
格子の一様化元 $t$ の冪による離散付値が入り、剰余は特殊隅の商可除代数となる。
完備化後も次数4は変わらず、中心の剰余が剰余可除代数の全中心である。
この付値で被約traceは整で剰余の被約traceと可換である。
実際、剰余の中心上の基底十六元を付値環の単元に持ち上げれば、完備中心の付値環上の
可除代数の付値環の基底になる。独立性は最小付値で剰余を取り、生成性は精度ごとの引き算と
完備性による。この基底での左正則作用のtraceを4で割ればよい。

一般点で各斉次経路 $b$ の被約traceは同次数の隅の中心に入る
（既知の正規中心上の有限な隅で、一般点では隅中心と全中心は一致する）。
格子の元なら付値が非負である。固定次数での付値非負は元の自由格子に属することそのものなので、
このtraceも同じ次数の格子に入る。特殊化すると $A$ は全ての被約traceで閉じている。
Newton公式から各元は中心係数の被約特性多項式を満たす。
従って $A$ は中心上有限である。詳しくは $A$ は有限個の長さ4の語で生成されるPI代数なので、
Shirshovの高さ定理により有限個の語の冪の有界個数の積で張られる。
それら有限語の被約特性係数で生成した中心部分環上、その冪の指数も3以下にできる。
よって有限加群になり、全中心 $Z(A)$ も有限生成である。
さらに特殊隅の任意の斉次中心 $c$ は、任意の持ち上げの被約traceの四分の一の特殊化である。
従って隅の中心格子は全次数で特殊中心へ全射し、一般中心とHilbert関数が一致する。

各 $M_i=e_iBe_0$ は最初に頂点0へ達する経路で右 $A$ 加群として有限生成である。
$M_i$ は右 $A$ 上torsion-freeである。$ma=0,\ a\ne0$ なら全ての $b\in e_0Be_i$ で
$bma=0$、従って $bm=0$ である。すると左から十分長い経路は必ず頂点0を通るため
$m$ を消す。AS性から有限次元の左片側イデアルはない（あればその次数付きsocleに
頂点単純が入り、次数零のAS消滅に反する）。各主張はまず斉次成分または
次数フィルトレーションで示せ、非斉次の場合も頂点単純だけを持つ有界な経路作用に同じ議論が使える。
よって $m=0$。$M_i$ と $A$ のHilbert多項式の最高係数は同じで、
同じ有限生成次数付き整中心上の加群の重複度比はその一般階数比なので、
$M_i$ は一般可除代数上階数1である。左乗法による $B\to\operatorname{End}_A(\bigoplus M_i)$ は忠実である。
核は右から十分長い全経路に消され、右側にも有限次元イデアルはないからである。
これで全隅の整性と全ての非零合成の非零性が従う。

一般の一つの中心の四成分の付値は全て一致する。
隅の積の非零性により隣接位置間の非零な経路を掛けても格子の付値は変わらないからである。
従って零番の中心格子の各元は全中心格子へ持ち上がり、その剰余は全位置で可換なままである。
他方零番への全中心の射は上の忠実性と非零合成で単射である。
よって $Z(B)=Z(A)$（零番制限で同定）であり、各 $M_i$ がこの上有限であることと
忠実な左乗法表示から $B$ も中心有限となる。
\end{proof}

\begin{proof}[四十八次における基点の排除]
中心錐の頂点以外の任意の閉点で単純表現 $S$ を一つ選ぶ。全頂点次元は同じ
$d\le4$ である。中心が今はNoetherで有限なため、次数を展開した加群
$M_n=Se_{n\bmod4}$ は有限生成で有限射影分解を持つ。
実際非零中心値を持つ正次数元が周期的同型を与える。従って
$(\sum_{i=0}^3\dim(Se_i)t^i)/(1-t^4)=p(t)/(1-t)^4$ となり四数が等しい。
上界はPI次数16による。各始点で、長さ4の経路のある線型結合が非冪零に作用することを示す。
$d=1$ は各辺非零なので明らかである。$d>1$ では各辺の各非零ベクトル上で
全矢の像が少なくとも二次元であり、双対でも同じである。
共通核を単純性で除き、像が直線なら有限点グラフで四段後も同じ四点を回り、
固有直線から薄い部分表現が出る定理\ref{thm:finite-two-dimensional-quadrics}の証明は
任意の頂点次元でそのまま働く。そのため任意の超平面は次の全空間へ送られる
（真部分なら消す一共変元の双対像が一次元以下になる）。

$d=2,3$ では従って任意の非零ベクトルが二段以内で次の全空間を生成する。
$d=4$ でもどの始点にも一段で少なくとも三次元を生成するベクトルがある。
なければ全ての $[v]\in\mathbb P^3$ で一定階数2となり、評価写像
$4\mathcal O(-1)\to4\mathcal O$ の像束 $E$ の行列式は $\mathcal O(-j)$、
$0\le j\le2$ である（$E^*$ と $E(1)$ は大域生成）。
$j=0$ なら $E$ は一定二平面となり全矢像が真部分である。
$j=1$ ならPl\"ucker写像は基点のない一次式で与えられ、三次元射影空間を
$\mathrm{Gr}(2,4)$ の中に線型に埋めるが、滑らかな四次元二次曲面にはそのような線型空間はない。
$j=2$ なら $E(1)$ は行列式自明な大域生成束で自明となる。
$2\mathcal O(-1)\hookrightarrow4\mathcal O$ の二列を二つの $v$ の線型写像と読むと、
どの非零な列の一次結合も全ての $[v]$ で非零で、可逆な四行四列行列を与える。
このような行列鉛筆は行列式の零点を持ち矛盾する。
以上よりあるベクトル $v$ が二段で全空間を生成し、その先も全空間を生成する。
四段で $v$ 自身を像に取る経路線型結合 $b$ が存在し、$b$ は非冪零である。

この $b$ のモニック被約特性多項式の四係数は、始点の隅から全中心に延びる
換算次数 $1,2,3,4$ の元である。全てがこの指標で零なら $b^4=0$ と作用して矛盾する。
$12$ はこの四次数全ての倍数なので、換算次数12の中心に頂点外の共通零点はない。
従ってその一般四元で切れば中心剰余は有限次元である。理由は
三次元射影中心を基点無しの豊富なVeronese次数の全系で有限射に移し、
四つの一般超平面との共通交叉を取ればよい。$B$ も中心上有限なのでその剰余も有限となる。
経路長48を定理\ref{thm:successively-central-cutoff}に代入して主張の有限判定を得る。
逆に検査を満たせばAS正則となり四中心上自由である。
パラメータの次数に関係なく、定理\ref{thm:central-cut-center-dichotomy}の
頂点射影の整数重複度1による証明で、合成可能な斉次非零元の積は非零となる。
\end{proof}
